% Part: first-order-logic
% Chapter: tableaux
% Section: rules-and-proofs

\documentclass[../../../include/open-logic-section]{subfiles}

\begin{document}

\iftag{FOL}
      {\olfileid[hi]{fol}{tab}{rul}}
      {\olfileid[hi]{pl}{tab}{rul}}

\olsection{नियम और \usetoken{P}{tableau}}

!!^a{tableau} इस बात का क्रमबद्ध सर्वेक्षण है कि !!a{sentence} किसी
!!a{structure} में किन संभावित तरीकों से सत्य या असत्य हो सकता है। यह
!!{signed formula}s से निर्मित होता है: !!{sentence}s के साथ सत्य-मूल्य का एक
``चिह्न'', जो या तो $\True$ होता है या~$\False$। इन चिह्नित !!{formula}s को
एक (नीचे की ओर बढ़ने वाले) वृक्ष में व्यवस्थित किया जाता है।

\begin{defn}
  \emph{!!{signed formula}} सत्य-मूल्य और !!a{sentence} से बना युग्म है,
  अर्थात् इन दोनों रूपों में से कोई एक:
  \[
  \sFmla{\True}{!A} \text{ अथवा } \sFmla{\False}{!A}.
  \]
\end{defn}

सहज रूप से, हम $\sFmla{\True}{!A}$ को ``$!A$ सत्य हो सकता है'' और
$\sFmla{\False}{!A}$ को ``$!A$ असत्य हो सकता है'' (किसी !!{structure} में)
पढ़ सकते हैं।

वृक्ष का प्रत्येक !!{signed formula} या तो \emph{मान्यता} है (मान्यताएँ वृक्ष
के बिल्कुल ऊपर लिखी जाती हैं), या वह अपने ऊपर के किसी !!{signed formula} से
किसी अनुमान-नियम द्वारा प्राप्त होता है। !!{main operator} का प्रत्येक संभावित
रूप—अर्थात् पूर्ववर्ती !!{formula} का प्रधान संक्रियक—दो नियमों से संबद्ध है: एक उस दशा के लिए
जिसमें चिह्न $\True$ है, और दूसरा उस दशा के लिए जिसमें चिह्न~$\False$ है। कुछ नियम वृक्ष
में शाखाएँ बनाने देते हैं, और कुछ केवल शाखा में !!{signed formula}s को जोड़ते
हैं। कोई नियम ठीक पहले आए !!{signed formula} पर ही नहीं, बल्कि मूल से लेकर
नियम लगाए जाने के स्थान तक की शाखा में स्थित किसी भी !!{signed formula} पर
लगाया जा सकता है (और प्रायः लगाना ही पड़ता है)।

कोई शाखा तब \emph{बंद} होती है जब उसमें $\sFmla{\True}{!A}$ और
$\sFmla{\False}{!A}$ दोनों हों। वह !!{tableau} बंद है जिसकी प्रत्येक शाखा
बंद है। सहज व्याख्या में हर शाखा एक संयुक्त संभावना बताती है, किन्तु
$\sFmla{\True}{!A}$ और $\sFmla{\False}{!A}$ का एक साथ संभव होना असंभव है।
दूसरे शब्दों में, किसी शाखा के बंद होने पर उसके द्वारा बताई गई संभावना निरस्त
हो जाती है। विशेषतः इसका अर्थ है कि बंद !!{tableau}, $\sFmla{\True}{!A}$ रूप
की प्रत्येक मान्यता को सत्य और $\sFmla{\False}{!A}$ रूप की प्रत्येक मान्यता
को असत्य बनाने की सभी संयुक्त संभावनाएँ निरस्त कर देता है।

$!A$ \emph{के लिए बंद !!{tableau}} ऐसा बंद !!{tableau} है जिसका
मूल~$\sFmla{\False}{!A}$ हो। यदि ऐसा बंद !!{tableau} अस्तित्व में है, तो
$!A$ के असत्य होने की सभी संभावनाएँ निरस्त हो चुकी हैं; अर्थात् $!A$ प्रत्येक
!!{structure} में सत्य होना चाहिए।

\end{document}
